\DocumentMetadata{
  pdfversion=2.0,
  pdfstandard=ua-2,
  lang=en-US,
  tagging=on,
  tagging-setup={math/setup=mathml-SE,tabsorder=structure}
}
\documentclass[11pt,letterpaper]{article}
\usepackage[margin=0.68in,includefoot]{geometry}
\usepackage{newcomputermodern}
\usepackage{amsmath,amscd,mathtools}
\usepackage{tikz,adjustbox}
\providecommand{\square}{\mdlgwhtsquare}
\usepackage[hidelinks]{hyperref}
\hypersetup{
  pdftitle={Algebra PhD Exam, September 2004},
  pdfauthor={University of Florida Department of Mathematics},
  pdfsubject={Graduate mathematics examination},
  pdfdisplaydoctitle=true,
  bookmarksopen=true
}
\setcounter{secnumdepth}{0}
\setlength{\parindent}{0pt}
\setlength{\parskip}{0.28em}
\setlength{\emergencystretch}{3em}
\setlength{\leftmargini}{2.15em}
\setlength{\leftmarginii}{2.65em}
\setlength{\leftmarginiii}{3em}
\renewcommand{\labelenumi}{\textbf{\theenumi.}}
\pagestyle{empty}
\raggedbottom
\allowdisplaybreaks
\newenvironment{examproblems}[1]{%
  \begin{enumerate}%
  \setcounter{enumi}{#1}%
  \setlength{\topsep}{0.35em}%
  \setlength{\partopsep}{0pt}%
  \setlength{\itemsep}{0.48em}%
  \setlength{\parsep}{0.18em}%
}{\end{enumerate}}
\newenvironment{examsubparts}{%
  \begin{enumerate}%
  \setlength{\topsep}{0.18em}%
  \setlength{\partopsep}{0pt}%
  \setlength{\itemsep}{0.16em}%
  \setlength{\parsep}{0.1em}%
}{\end{enumerate}}
\newenvironment{examinstructions}{%
  \par\begingroup\itshape
}{%
  \par\endgroup\vspace{0.18em}
}
\makeatletter
\renewcommand\section{\@startsection{section}{1}{0pt}%
  {-1.25ex plus -.25ex minus -.1ex}%
  {0.55ex plus .1ex}%
  {\normalfont\large\bfseries}}
\renewcommand{\@maketitle}{%
  \newpage\null\vskip -1em%
  {\footnotesize\bfseries
    \href{https://www.ufl.edu/}{University of Florida}, \href{https://math.ufl.edu/}{Department of Mathematics}\par}%
  \vskip -0.5em%
  \tagstructbegin{tag=Title}%
    {\LARGE\bfseries\href{https://gma.math.ufl.edu/exams/algebra-phd/}{Algebra PhD Exam}, September 2004\par}%
  \tagstructend%
  \vskip 0.25em%
  \tagmcbegin{artifact}%
    \hrule height 1pt%
  \tagmcend%
  \vskip 1em%
}
\makeatother
\title{Algebra PhD Exam, September 2004}
\author{}
\date{}
\begin{document}
\maketitle
\thispagestyle{empty}
\begin{examinstructions}
Answer seven problems. Write your answers clearly in complete English sentences. You may quote results (within reason) as long as you state them clearly.
\end{examinstructions}
\begin{examproblems}{0}
\item
This problem has four parts.
\begin{examsubparts}
\item[\textnormal{(a)}]
Define concrete category.
\item[\textnormal{(b)}]
Define what it means for an object~\(F\) in a concrete category to be free on a set~\(S\).
\item[\textnormal{(c)}]
Given a set~\(S\), construct a free object on~\(S\) in the category of abelian groups.
\item[\textnormal{(d)}]
Given a set~\(S\), construct a free object on~\(S\) in the category of sets.
\end{examsubparts}
\item
Let~\(F\) be a field, let~\(A\) be a central simple algebra over~\(F\), and let~\(B\) be a simple~\(F\)-algebra with 1. Prove that \(A\otimes_F B\) is simple.
\item
Use the Wedderburn–Artin Theorem to determine all isomorphism classes of semisimple rings of order \(720=2^4\cdot 3^2\cdot 5\).
\item
Let \(\zeta=e^{2\pi i/16}\) be a primitive 16th root of unity and let \(F=\mathbb{Q}(\zeta)\). Determine the Galois group \(G=\operatorname{Gal}(F/\mathbb{Q})\), and for each subgroup \(H\leq G\) determine the fixed field \(F^H\).
\item
Let~\(p\) be prime and let~\(n\) be a positive integer. Prove that there is a field \(\mathbb{F}_{p^n}\) with \(p^n\) elements, and that this field is uniquely determined up to isomorphism.
\item
Let \(G=\langle x,y:x^3=y^2,\ y^4=1\rangle\). Prove that~\(G\) is infinite and non-abelian.
\item
Give an example of a projective module which is not free. Justify your answer.
\item
Let~\(R\) be an Artinian commutative ring with 1. Prove that~\(R\) has only finitely many prime ideals, all of which are maximal.
\item
Let~\(R\) be a commutative ring with 1 and let~\(S\) be a multiplicative subset of~\(R\) which contains no zero divisors. Prove that there is a one-to-one correspondence between prime ideals of \(S^{-1}R\) and prime ideals~\(P\) of~\(R\) such that \(P\cap S=\varnothing\).
\item
Let~\(R\) be a commutative ring with 1.
\begin{examsubparts}
\item[\textnormal{(a)}]
Prove that if \(I_1\) and \(I_2\) are ideals in~\(R\) such that \(I_1\cap I_2=P\) is prime then either \(I_1=P\) or \(I_2=P\).
\item[\textnormal{(b)}]
Is the statement above still true if “prime” is replaced by “primary”? Justify your answer.
\end{examsubparts}
\item
Let \(R\subset S\) be commutative rings with \(1_R=1_S\) such that~\(S\) is integral over~\(R\). Let~\(C\) be an algebraically closed field and let \(\phi:R\to C\) be a unitary ring homomorphism. Prove that there is a ring homomorphism \(\widetilde{\phi}:S\to C\) such that \(\widetilde{\phi}|_R=\phi\).
\end{examproblems}
\end{document}
